\chapter{The $K3 \times E$ Tower and the Igusa Cusp Form}
\label{ch:k3-times-e}

This chapter develops the prototypical example of a quantum vertex chiral group: the generalized BKM superalgebra $\mathfrak{g}_{\Delta_5}$ attached to the CY3 family $X = (S \times E)/(\mathbb{Z}/N\mathbb{Z})$, whose denominator identity is the Igusa cusp form $\Delta_5$.

\section{The CY3 geometry}
\label{sec:k3e-geometry}

Let $(E, e_0)$ be an elliptic curve with an $N$-torsion point and $S$ a K3 surface with elliptic fibration $\pi \colon S \to \mathbb{P}^1$ admitting sections $s_1, s_2 \colon \mathbb{P}^1 \to S$ with $s_2$ of order $N$ relative to $s_1$.  The product $S \times E$ admits a free $\mathbb{Z}/N\mathbb{Z}$-action
\[
  (s, e) \longmapsto (s + s_2(\pi(s)), e + e_0),
\]
and the quotient $X = (S \times E)/(\mathbb{Z}/N\mathbb{Z})$ is a projective Calabi--Yau threefold.

\begin{definition}[The DT zeta function]
\label{def:dt-zeta-k3e}
Let $F$ be the fiber class of $\pi$ and $\beta_h = \frac{1}{N}(s_1 + \cdots + s_N + hF)$ for $h \geq 0$.  The DT zeta function of $X$ is
\[
  Z^X(q, t, p) = \sum_{h=0}^{\infty} \sum_{d=0}^{\infty} \sum_{n \in \mathbb{Z}} \mathrm{DT}^X_{n,(\beta_h, d)} \, q^{d-1} \, t^{\frac{1}{2}\langle \beta_h, \beta_k \rangle} (-p)^n.
\]
\end{definition}

\begin{theorem}[Oberdieck--Pixton]
\label{thm:dt-igusa}
For $X = S \times E$ of order $N = 1$:
\[
  Z^X(q, t, p) = \frac{C}{(\Delta_5)^2}
\]
for a constant $C$, where $\Delta_5$ is the Igusa cusp form of weight $5$ for $\mathrm{Sp}_4(\mathbb{Z})$.
\end{theorem}

The elliptic genus of K3 --- the weak Jacobi form $\phi_{0,1}$ of weight $0$ and index $1$ --- governs the root multiplicities of the BKM superalgebra.

\section{The lattice $\Lambda^{3,2}$ and the isomorphism $\mathrm{Sp}_4 \simeq \mathrm{SO}_+$}
\label{sec:k3e-lattice}

Consider the rank-4 free $\mathbb{Z}$-module $\Lambda^4 = \mathbb{Z}e_1 \oplus \mathbb{Z}e_2 \oplus \mathbb{Z}e_3 \oplus \mathbb{Z}e_4$.  The exterior square $\bigwedge^2 \colon \mathrm{SL}_4(\mathbb{Z}) \to \mathrm{GL}(\bigwedge^2 \Lambda^4)$ and the Pfaffian scalar product $(u,v) = u \wedge v / (e_1 \wedge e_2 \wedge e_3 \wedge e_4)$ give a signature $(3,3)$ form.  The lattice
\[
  \Lambda^{3,2} = (e_1 \wedge e_3 + e_2 \wedge e_4)^\perp \subset \bigwedge^2 \Lambda^4
\]
has signature $(3,2)$ and satisfies $\Lambda^{3,2} \simeq \Lambda^{1,1} \oplus \Lambda^{1,1} \oplus [2]$.

\begin{lemma}
\label{lem:sp4-so-iso}
The correspondence $\bigwedge^2$ defines an isomorphism
\[
  % B-class fix: Sp_4 is 4x4, so quotient by +-I_4; O(Lambda^{3,2}) is
  % 5x5, so quotient by +-I_5.
  \bigwedge^2 \colon \mathrm{Sp}_4(\mathbb{Z})/\{\pm I_4\} \xrightarrow{\;\sim\;} \mathrm{O}(\Lambda^{3,2})_+/\{\pm I_5\}.
\]
\end{lemma}

In the basis $(f_1, f_2, f_3, f_{-2}, f_{-1})$ with $f_1 = e_1 \wedge e_2$, $f_2 = e_2 \wedge e_3$, $f_3 = e_1 \wedge e_3 - e_2 \wedge e_4$, $f_{-2} = e_4 \wedge e_1$, $f_{-1} = e_4 \wedge e_3$, the orthogonal group $\mathrm{O}_\mathbb{R}(\Lambda^{3,2})_+$ acts on the type IV domain $\mathbb{H}^{IV}_+$, which identifies with the Siegel upper half-space $\mathbb{H}_2$.

\section{The hyperbolic sublattice and reflection groups}
\label{sec:k3e-reflections}

Fix the primitive hyperbolic sublattice
\[
  \Lambda^{2,1} = \Lambda^{1,1} \oplus [2] \simeq \mathbb{Z}f_2 \oplus \mathbb{Z}f_3 \oplus \mathbb{Z}f_{-2} \subset \Lambda^{3,2}.
\]
The sublattice $\Lambda^{2,1}_{II}$ generated by elements of type II (those $\delta$ with $(\delta, \Lambda^{2,1}) = 2\mathbb{Z}$) has three real simple roots:
\[
  \delta_1 = 2f_2 - f_3, \qquad \delta_2 = 2f_{-2} - f_3, \qquad \delta_3 = f_3,
\]
with Gram matrix
\[
  (\delta_i, \delta_j) = \begin{pmatrix} 2 & -2 & -2 \\ -2 & 2 & -2 \\ -2 & -2 & 2 \end{pmatrix}.
\]

The fundamental polyhedron $\mathcal{P}_{II}$ has $\mathrm{Aut}(\mathcal{P}_{II}) \simeq S_3$, and
\[
  \mathrm{O}(\Lambda^{2,1})_+ \simeq W^{(2)}(\Lambda^{2,1}_{II}) \rtimes \mathrm{Aut}(\mathcal{P}_{II}).
\]

\begin{definition}[Weyl vector]
\label{def:k3e-weyl-vector}
The Weyl vector $\rho \in \mathcal{P}_{II} \otimes \mathbb{Q}$ satisfies $(\rho, \delta_i) = -(\delta_i, \delta_i)/2 = -1$ for all $i$:
\[
  \rho = \tfrac{1}{2}\delta_1 + \tfrac{1}{2}\delta_2 + \tfrac{1}{2}\delta_3 = f_2 - \tfrac{1}{2}f_3 + f_{-2}.
\]
\end{definition}

\section{The generalized BKM superalgebra $\mathfrak{g}_{\Delta_5}$}
\label{sec:k3e-bkm}

The automorphic correction of the Kac--Moody algebra $\mathfrak{g}$ (with Gram matrix $(\delta_i, \delta_j)$) by the Fourier coefficients of $\Delta_5$ produces the generalized BKM Lie superalgebra $\mathfrak{g}_{\Delta_5}$.

\begin{construction}[Root system of $\mathfrak{g}_{\Delta_5}$]
\label{constr:k3e-roots}
The root system decomposes as $\Delta = \Delta^{\mathrm{re}} \sqcup \Delta^{\mathrm{im}}_0 \sqcup \Delta^{\mathrm{im}}_1$:
\begin{itemize}
  \item \textbf{Real even simple roots}: $\Delta^{\mathrm{re}}_0 = \Delta^{\mathrm{re}} = \mathcal{P}_{II,\mathrm{prim}} = \{\delta_1, \delta_2, \delta_3\}$.
  \item \textbf{Even imaginary simple roots}: $\Delta^{\mathrm{im}}_0 = \{ \tau(a) \cdot a \mid (a,a) = 0, \, \tau(a) > 0 \}$, repeated $\tau(a)$ times.
  \item \textbf{Odd imaginary simple roots}: $\Delta^{\mathrm{im}}_1 = \{ m(a) \cdot a \mid (a,a) < 0, \, m(a) < 0 \}$, where the element $a$ is repeated $-m(a)$ times and these generators are fermionic.
\end{itemize}
Here $m(a) = -\frac{1}{64} f(n,l,m)$ for $a \in \Lambda^{2,1}_{II} \cap \mathbb{R}_{\geq 0} \mathcal{P}_{II}$ corresponding to the Fourier coefficient $f(n,l,m)$ of $\Delta_5$.
\end{construction}

The superalgebra $\mathfrak{g}_{\Delta_5}$ has the triangular decomposition
\[
  \mathfrak{g}_{\Delta_5} = \biggl(\bigoplus_{\alpha \in \widetilde{C(\Lambda^{2,1}_{II})_+}} \mathfrak{g}_\alpha\biggr) \;\oplus\; (\Lambda^{2,1}_{II} \otimes \mathbb{R}) \;\oplus\; \biggl(\bigoplus_{\alpha \in -\widetilde{C(\Lambda^{2,1}_{II})_+}} \mathfrak{g}_\alpha\biggr)
\]
with generators $h_\alpha, e_\alpha, f_\alpha$ for $\alpha \in \Delta$ satisfying the standard BKM relations.

\section{The denominator identity: $\Delta_5 = \Phi$}
\label{sec:k3e-denominator}

\begin{theorem}[Denominator identity]
\label{thm:k3e-denominator}
The Weyl--Kac--Borcherds character formula applied to the trivial one-dimensional representation of $\mathfrak{g}_{\Delta_5}$ gives
\[
  \frac{1}{64} \Delta_5(2Z) = \Phi(z)
\]
where $\Phi$ is the denominator function
\[
  \Phi = \sum_{w \in W^{(2)}(\Lambda^{2,1})} \det(w) \exp(-\pi i \langle w(\rho), z \rangle) - \sum_{a \in \Lambda^{2,1}_{II} \cap \mathbb{R}_{\geq 0} \mathcal{P}_{II}} m(a) \exp(-\pi i \langle w(\rho + a), z \rangle).
\]
Equivalently, in product form:
\[
  \Phi = \exp(-2\pi i \langle \rho, z \rangle) \prod_{\alpha \in \Delta_+} (1 - \exp(-2\pi i \langle \alpha, z \rangle))^{\mathrm{mult}\, \alpha}.
\]
\end{theorem}

\begin{theorem}[Product formula for $\Delta_5$]
\label{thm:k3e-product}
\[
  \frac{1}{64} \Delta_5(Z) = -\exp(\pi i(z_1 + z_2 + z_3)) \prod_{\substack{n,l,m \in \mathbb{Z}, \, (n,l,m) > 0}} (1 - \exp(2\pi i(nz_1 + lz_2 + mz_3)))^{f(nm,l)}
\]
where $Z = \bigl(\begin{smallmatrix} z_1 & z_2 \\ z_2 & z_3 \end{smallmatrix}\bigr) \in \mathbb{H}_2$ and $f(n,l)$ are the Fourier coefficients of the weak Jacobi form $\phi_{0,1}$ of weight $0$ and index $1$.
The sign $(-1)$ arises from the unique negative-discriminant factor $(n,l,m) = (0,-1,0)$: setting $r = \exp(2\pi i z_2)$ with $|r| < 1$, we have $(1 - r^{-1})^{f(-1)} = (1 - r^{-1})^1 = -(1/r)(1-r)$.
Absorbing this factor, the sign-free form reads
\[
  \frac{1}{64} \Delta_5(Z) = \exp(\pi i(z_1 - z_2 + z_3)) \cdot (1 - \exp(2\pi i z_2)) \prod_{\substack{(n,l,m) > 0 \\ (n,l,m) \neq (0,-1,0)}} (1 - \exp(2\pi i(nz_1 + lz_2 + mz_3)))^{f(nm,l)}.
\]
\ClaimStatusProvedHere
\end{theorem}

\section{The weak Jacobi form $\phi_{0,1}$ and root multiplicities}
\label{sec:k3e-phi01}

The weak Jacobi form $\phi_{0,1} = \phi_{12,1}/\delta_{12}$ (where $\delta_{12} = q \prod_{n \geq 1} (1 - q^n)^{24}$ is the weight-12 cusp form) is the K3 elliptic genus.  Its Fourier coefficients $f(n,l)$ are the super-dimensions of the root spaces of $\mathfrak{g}_{\Delta_5}$:
\[
  \mathrm{mult}\, \alpha = f(nm, l) \quad \text{for } \alpha = (n,l,m) \in \Delta_+.
\]

\section{The quantum vertex chiral group $G(K3 \times E)$}
\label{sec:k3e-qvcg}

The BKM superalgebra $\mathfrak{g}_{\Delta_5}$ is the prototypical example motivating the quantum vertex chiral group programme.  In the language of Volumes~I--III:

\begin{itemize}
  \item \textbf{(Theorem.)} The generalized root datum $\mathcal{R}(K3 \times E)$ is $(\Lambda^{3,2}, \Delta^{\mathrm{re}}, \Delta^{\mathrm{im}}_0, \Delta^{\mathrm{im}}_1, W^{(2)}(\Lambda^{2,1}_{II}), \rho, f(nm,l))$.  This is a mathematical fact (Gritsenko--Nikulin).
  \item \textbf{(Conjecture.)} There exists a chiral algebra $A_{K3 \times E} = \Phi(D^b(\mathrm{Coh}(X)))$ whose bar complex $B(A)$ encodes the product formula for $\Delta_5$.  \emph{Note}: the functor $\Phi$ is constructed for $d = 2$ (Theorem~CY-A$_2$); the $d = 3$ case (which applies here, since $K3 \times E$ is CY$_3$) is the content of Conjecture~CY-A$_3$.
  \item \textbf{(Observation/Conjecture.)} The number $5 = \mathrm{wt}(\Delta_5) = h^{1,1}(K3)/4$ appears in the structural position of a modular characteristic: $\kappa = 5$.  Without the chiral algebra $A_{K3 \times E}$, this identification is an observation matching the Borcherds lift weight formula, not a computation from the Volume~I definition of $\kappa$.
  % AP42: This identification holds at the level of formal Euler products
  % and Fourier coefficients, but the naive BCH pair-commutator does NOT
  % reproduce phi_{0,1} multiplicities (quantitative mismatch). The gap
  % measures higher BPS bound-state contributions requiring full motivic DT
  % theory. State the level of validity explicitly.
  \item \textbf{(Theorem at the formal product level.)} The automorphic correction $\mathfrak{g} \leadsto \mathfrak{g}_{\Delta_5}$ matches the shadow obstruction tower structure at the level of the formal Euler product: the arity-$r$ projection captures imaginary roots of BPS charge $\leq r$.  This is computationally verified at arities 2--6.  \emph{Caveat}: the naive BCH pair-commutator does not reproduce $\phi_{0,1}$ multiplicities at low arities; the full identification requires the motivic Hall algebra level (AP42).
  \item \textbf{(Conjectural.)} The $E_2$-structure should come from the $\mathrm{Sp}_4(\mathbb{Z})$-action on $\mathbb{H}_2$, connecting genus-2 modular structure to braided monoidal structure via the modular functor.  Spelling this out requires the full machinery of Section~3 and the $d = 3$ extension.
\end{itemize}

\begin{conjecture}[Eight quantum vertex chiral groups]
\label{conj:eight-qvcg}
Each of the eight diagonal divisor Siegel modular forms arises as the denominator identity of a quantum vertex chiral group $G(X_N)$ attached to the CY3 $X_N = (S \times E)/(\mathbb{Z}/N\mathbb{Z})$, with root multiplicities from the $g_N, h_M$-twisted twined elliptic genera of K3.
\end{conjecture}
